\section{Задание}

\textbf{Вариант: } \\
	1. Организовать хранение коллекции в реляционной СУБД (PostgresQL). Убрать хранение коллекции в файле.\\
	2. Для генерации поля id использовать средства базы данных (sequence).\\
	3. Обновлять состояние коллекции в памяти только при успешном добавлении объекта в БД\\
	4. Все команды получения данных должны работать с коллекцией в памяти, а не в БД\\
	5. Организовать возможность регистрации и авторизации пользователей. У пользователя есть возможность указать пароль.\\
	6. Пароли при хранении хэшировать алгоритмом SHA-256\\
	7. Запретить выполнение команд не авторизованным пользователям.\\
	8. При хранении объектов сохранять информацию о пользователе, который создал этот объект.\\
	9. Пользователи должны иметь возможность просмотра всех объектов коллекции, но модифицировать могут только принадлежащие им.\\
	10. Для идентификации пользователя отправлять логин и пароль с каждым запросом.\\

\textbf{Необходимо реализовать многопоточную обработку запросов.}\\
	1. Для многопоточного чтения запросов использовать\\
	2. Для многопотчной обработки полученного запроса использовать .\\
	3. Для многопоточной отправки ответа использовать\\
	4. Для синхронизации доступа к коллекции использовать\\
	

\textbf{Порядок выполнения работы::}\\
	1.В качестве базы данных использовать PostgreSQL.\\
	2. Для подключения к БД на кафедральном сервере использовать хост pg, имя базы данных - studs, имя пользователя/пароль совпадают с таковыми для подключения к серверу.\\
\textbf{Отчёт по работе должен содержать:} \\
	1. Текст задания.\\
	2. Диаграмма классов разработанной программы.\\
	4. Исходный код программы.\\
	5. Выводы по работе.\\
	
\textbf{Вопросы к защите лабораторной работы:}\\
	1. Многопоточность. Класс Thread, интерфейс Runnable. Модификатор synchronized.\\
	2. Методы wait(), notify() класса Object, интерфейсы Lock и Condition\\
	3. Классы-сихронизаторы из пакета java.util.concurrent.\\
	4. Модификатор volatile. Атомарные типы данных и операции.\\
	5. Коллекции из пакета java.util.concurrent.\\
	6. Интерфейсы Executor, ExecutorService, Callable, Future.\\
	7. Пулы потоков.\\
	8. JDBC. Порядок взаимодействия с базой данных. Класс DriverManager. Интерфейс Connection.\\
	9. Интерфейсы Statement, PreparedStatement, ResultSet, RowSet.\\
	10. Шаблоны проектирования.\\

\newpage
\section{Код и Диаграмм}
\textbf{Исходный код доступен по ссылке или QR-коду:}\\
\\     
\underline{$https://github.com/ndwannafly/Programming-Lab-2nd-Semester/tree/main/LAB_7$}\\


\begin{figure}[H]
\includegraphics[scale=0.8]{img/SourceCode}
\label{pic:SourceCode}
\end{figure}

\underline{$https://github.com/ndwannafly/Programming-Lab-2nd-Semester/blob/main/LAB_7/report$}\\


\begin{figure}[H]
\includegraphics[scale=0.8]{img/QRDiagram}
\label{pic:QRDiagram}
\end{figure}

\section{Вывод}
Во время выполнения лабораторной работы мы изучили многопоточное программирование, как работать с базой данных (PostgresQL), алгоритм хеширования и организовали регистрацию, вход, выход для пользователей.